\chapter{Semantics Before Structures}
\label{ch:pebble-semantics}

Two implementations can agree on the fields of an addition instruction and disagree on its result.
Before choosing a container, specify the machine that the container describes.
PebbleIR has enough control flow to expose a subtler problem too: returning the right number is not the same as preserving whether the program returns.

\section{Values and observations}
\label{sec:pebble-observations}
The type \code{u8} denotes $\{0,\ldots,255\}$.
The type \code{i1} denotes $\{0,1\}$, with zero false and one true.
The instruction \code{lt} compares two unsigned byte values and returns an \code{i1}.
There is no implicit conversion from a byte to a branch condition.
The verifier rejects branching directly on \code{x}, even when a particular test input happens to be zero or one.

For a defined input, a valid function either returns a byte or continues executing forever.
These are its semantic observations.
\term{Fuel exhaustion} is a third result of our testing interface, not an additional instruction in the language and not proof of divergence.
An invalid representation is rejected before execution and therefore is not a program whose behavior an optimizer must preserve.

\section{A state that explains a step}
\label{sec:pebble-state}
An interpreter state contains the current block, the predecessor just taken, and a register map from value IDs to bytes or booleans.
Entering a block first selects every phi input from the predecessor and captures all those values.
It then assigns the phi results simultaneously, executes ordinary instructions in ID order within the block, and executes the terminator.
Figure~\ref{fig:pebble-step} separates these events.

\begin{figure}[htbp]
\centering
\begin{tikzpicture}[node distance=9mm]
\node[basic block] (select) {Select incoming phi values};
\node[basic block,below=of select] (assign) {Publish all phi results};
\node[basic block,below=of assign] (body) {Execute ordinary instructions};
\node[basic block,below=of body] (term) {Return or choose successor};
\draw[cfg edge] (select)--(assign);
\draw[cfg edge] (assign)--(body);
\draw[cfg edge] (body)--(term);
\end{tikzpicture}
\caption{One block entry selects a complete incoming state before executing the body.}
\label{fig:pebble-step}
\end{figure}

Consider a loop that swaps two values on its back edge.
If their old values are 1 and 2, sequentially assigning \code{a = b} and then \code{b = a} produces 2 and 2.
Parallel selection produces 2 and 1.
After a second iteration it produces 1 and 2 again.
The executable swap fixture checks precisely this distinction.

\begin{algorithm}[htbp]
\caption{Evaluating a block's phi prefix.}
\label{alg:pebble-phi-state}
\begin{pseudocode}
\Require verified block $B$, incoming predecessor $P$, register map $R$
\Ensure all phi results receive the values selected from the old $R$
\ForAll{phis $v=\phifn{p:a,q:b}$ in $B$}
\State $T[v]\gets R[a]$ if $P=p$, otherwise $R[b]$
\EndFor
\ForAll{phis $v$ in $B$}
\State $R[v]\gets T[v]$
\EndFor
\end{pseudocode}
\end{algorithm}

Algorithm~\ref{alg:pebble-phi-state} takes $O(P)$ time and scratch slots for $P$ phis when the prefix is directly indexed.
The teaching implementation scans the value table to find the prefix, adding $O(N)$ work per block entry.
Keeping that cost visible prevents a simple representation from acquiring an undeserved linear-execution claim.

\section{The host is not the language}
\label{sec:pebble-host}
The C implementation uses \code{unsigned} and masks arithmetic results with 255.
The largest sum is 510; the largest product is 65,025.
Both fit the minimum range required of C's unsigned integer type.
The comparison uses those already normalized unsigned values.
Consequently these operations implement the byte semantics without relying on signed overflow, optional eight-bit types, or a particular endianness.

Adding division would require a zero-divisor rule before adding an opcode.
Adding memory would require observations of stores and a model of invalid accesses.
Adding floating point would require an explicit decision about rounding and exceptional values.
The representation does not answer these questions merely because a host C operator exists.

\section{Exercises}
\label{sec:pebble-two-exercises}
\begin{enumerate}
\item \label{ex:pebble-swap} Trace two iterations of the swap using both sequential and parallel assignments. Identify the first incorrect state.
\item \label{ex:pebble-fuel} A program returns after 20 block entries. What can be concluded from exhaustion at budgets 5 and 19? What changes when it returns at budget 20?
\end{enumerate}

\section{Further reading}
\label{sec:pebble-two-reading}
The C11 draft specifies unsigned arithmetic and minimum integer ranges; use those requirements to justify the host implementation rather than observations on one machine~\cite[Sections~5.2.4.2.1 and~6.2.5]{wg14-2011-c11}.
